\documentclass[11pt]{article}
\usepackage{fontspec,amsmath,amssymb,amsthm,hyperref}
\usepackage[margin=1in]{geometry}
\setlength{\emergencystretch}{3em}
\title{Product and coordinate identities for the true Langevin generator}
\author{LangevinLean mathematical exposition}
\date{}
\begin{document}\maketitle
\noindent Source module: \texttt{AverageLMC/GeneratorProduct.lean}. Compile this proof document with XeLaTeX or LuaLaTeX. The accompanying HTML pairs the mathematical argument with the exact Lean statements and records the scope of the formalization.


The cutoff approximation for coordinate scores uses the exact generator product rule. These identities are deterministic calculus statements about the actual gradient, derivative and generator; no probabilistic cancellation is postulated.



For \(\mathcal L_f\psi=\Delta\psi-\langle\nabla f,\nabla\psi\rangle\), the compact cutoff construction uses the genuine product rule and the coordinate gradient identity:


\[\mathcal L_f(\chi\psi)=\chi\mathcal L_f\psi+\psi\mathcal L_f\chi+2\langle\nabla\chi,\nabla\psi\rangle.\]
\[\mathcal L_f(\partial_jf)=(G_f)_j,\qquad G_f=\nabla\Delta f-D(\nabla f)[\nabla f].\]
\section{Differentiate the product once and twice}


The Fréchet product rule, evaluated on a fixed direction, gives Dᵥ(χψ)=χDᵥψ+ψDᵥχ. The corresponding addition rule allows this expression to be differentiated once more. For C\ensuremath{{}^{2}} factors their first directional derivatives are differentiable, so another product rule gives Dᵥ\ensuremath{{}^{2}}(χψ)=χDᵥ\ensuremath{{}^{2}}ψ+ψDᵥ\ensuremath{{}^{2}}χ+2DᵥχDᵥψ. The cross coefficient is exactly two.



\noindent Lean declarations: \texttt{directionalDerivative\_mul}, \texttt{directionalDerivative\_add}, \texttt{directionalDerivative\_two\_mul}.

\section{Sum the directional generator product rules}


Subtract the first-order gradient drift from the second directional derivative product identity. Algebra separates the two factors’ directional generators and retains the cross term 2DᵥχDᵥψ. Summing over the standard orthonormal basis produces the full Langevin generator product rule. The sum of the products of the directional derivatives equals the inner product of the two actual gradients by the orthonormal-basis inner-product identity.



\noindent Lean declarations: \texttt{directionalGenerator\_mul}, \texttt{langevinGenerator\_mul}.

\section{Identify the coordinate score generator}


A coordinate score is the continuous linear projection of the C\ensuremath{{}^{2}} vector field ∇f. Projection commutes with its vector Laplacian, and the proved C\ensuremath{{}^{3}} commutation identity turns this into the corresponding coordinate of ∇Δf. The Fréchet chain rule identifies its derivative along ∇f with the same coordinate of D(∇f)∇f. Subtracting the drift term therefore gives exactly that coordinate of the actual vector gradientGenerator.



\noindent Lean declarations: \texttt{langevinGenerator\_coordinate\_gradient}.

\end{document}
